% Project Synopsis. Compile with pdfLaTeX (run twice so the contents page fills in).
% Edit names and PRNs in team.tex.
\documentclass[12pt,a4paper]{article}
% ============================================================================
% Team and institute details. EDIT THIS FILE ONLY; all three reports read it.
% Items marked FILL are left blank on purpose: they were not in the project
% files and must not be guessed.
% ============================================================================
\newcommand{\ProjectTitle}{Unified Content Authenticity Verification System}
\newcommand{\ProjectSubtitle}{Deepfake Detection with Cross-Generator Generalization and Document Tamper Localization}
\newcommand{\ProjectTitleCaps}{UNIFIED CONTENT AUTHENTICITY VERIFICATION SYSTEM: DEEPFAKE DETECTION WITH CROSS-GENERATOR GENERALIZATION AND DOCUMENT TAMPER LOCALIZATION}
\newcommand{\TeamName}{The Perceptrons}

\newcommand{\InstituteAddress}{GAT NO: 720/1 \& 2, NAGAR ROAD, WAGHOLI, PUNE-412207}
\newcommand{\AcademicYear}{2026--2027}
\newcommand{\AcademicYearShort}{2026-27}

% ---- students (upper-case on the cover, as in the department's format) ---------
\newcommand{\StudentOne}{PAWAN SINGH JAGDISH BHANDARI}
\newcommand{\StudentTwo}{RAHUL H. BRAMHANKAR}      % FILL: full name with father's name if the university needs it
\newcommand{\StudentThree}{BHOOMI KUSHWAH}          % FILL: full name with father's name if the university needs it
\newcommand{\StudentOneName}{Pawan Singh Jagdish Bhandari}
\newcommand{\StudentTwoName}{Rahul H. Bramhankar}
\newcommand{\StudentThreeName}{Bhoomi Kushwah}

% ---- PRN numbers: FILL (printed as a blank line until you type them) -----------
\newcommand{\PRNOne}{\underline{\hspace{3cm}}}
\newcommand{\PRNTwo}{\underline{\hspace{3cm}}}
\newcommand{\PRNThree}{\underline{\hspace{3cm}}}

% ---- guide, head, principal -------------------------------------------------------
\newcommand{\GuideName}{Dr. G. M. Bhandari}          % the research paper lists her as guide and she is also HOD
\newcommand{\GuideTitle}{Guide}
\newcommand{\HODName}{Dr. G. M. Bhandari}
\newcommand{\PrincipalName}{\underline{\hspace{3.6cm}}}   % FILL: current principal's name

% ============================================================================
% Shared preamble for the three project reports (Abstract, Synopsis, Stage 1).
% Every main.tex starts with \documentclass and then % ============================================================================
% Shared preamble for the three project reports (Abstract, Synopsis, Stage 1).
% Every main.tex starts with \documentclass and then % ============================================================================
% Shared preamble for the three project reports (Abstract, Synopsis, Stage 1).
% Every main.tex starts with \documentclass and then \input{preamble}.
% Only standard packages that are on Overleaf by default are used.
% ============================================================================
\usepackage[T1]{fontenc}
\usepackage[utf8]{inputenc}
\usepackage{mathptmx}                       % Times-style text and maths
\usepackage[a4paper,left=3cm,right=2.5cm,top=2.5cm,bottom=2.6cm]{geometry}
\usepackage{setspace}
\usepackage{graphicx}
\usepackage{xcolor}
\usepackage{booktabs,array,tabularx,multirow,longtable}
\usepackage{amsmath,amssymb}
\usepackage{enumitem}
\usepackage[font=small,labelfont=bf]{caption}
\usepackage{float}
\usepackage{fancyhdr}
\usepackage{titlesec}
\usepackage{url}
\usepackage[hidelinks]{hyperref}
\usepackage{tikz}
\usetikzlibrary{arrows.meta,positioning,shapes.geometric,shapes.misc,calc,fit,backgrounds,decorations.pathreplacing}
\usepackage{pgfplots}
\pgfplotsset{compat=1.16}
\usepackage{pgfgantt}

% ---- colours (also readable in black and white) ----------------------------
\definecolor{cBlue}{HTML}{DCE8F5}
\definecolor{cGreen}{HTML}{DDEEDD}
\definecolor{cOrange}{HTML}{FBE5CC}
\definecolor{cGrey}{HTML}{EEEEEE}
\definecolor{cLine}{HTML}{333333}

% ---- page style: footer as in the department's report format -----------------
\newcommand{\FooterText}{BSIOTR, Department of Computer Engineering, \AcademicYearShort}
\pagestyle{fancy}
\fancyhf{}
\renewcommand{\headrulewidth}{0pt}
\fancyfoot[C]{\small \FooterText\\[2pt]\thepage}
\fancypagestyle{plain}{\fancyhf{}\renewcommand{\headrulewidth}{0pt}\fancyfoot[C]{\small \FooterText\\[2pt]\thepage}}

% ---- headings ------------------------------------------------------------------
\makeatletter
\@ifundefined{chapter}{}{%
  \titleformat{\chapter}[display]{\centering\bfseries\Large}{CHAPTER \thechapter}{8pt}{\Large\MakeUppercase}
  \titlespacing*{\chapter}{0pt}{-10pt}{22pt}}
\makeatother
\titleformat{\section}{\bfseries\large}{\thesection}{0.6em}{}
\titleformat{\subsection}{\bfseries\normalsize}{\thesubsection}{0.6em}{}
\setlength{\parindent}{0pt}
\setlength{\parskip}{6pt}
\setlist{itemsep=2pt,topsep=3pt}
\renewcommand{\arraystretch}{1.18}
\setlength{\emergencystretch}{3em}

% ---- TikZ styles shared by all diagrams ----------------------------------------
\tikzset{
  >=Stealth,
  box/.style={rectangle,rounded corners=2pt,draw=cLine,thick,fill=cBlue,align=center,
              minimum height=9mm,inner sep=4pt,font=\small},
  boxg/.style={box,fill=cGreen},
  boxo/.style={box,fill=cOrange},
  boxw/.style={box,fill=white},
  designed/.style={box,fill=white,dashed},
  store/.style={rectangle,draw=cLine,thick,fill=cGrey,align=center,minimum height=8mm,inner sep=4pt,font=\small},
  actor/.style={rectangle,draw=cLine,thick,fill=cOrange,align=center,minimum height=8mm,inner sep=4pt,font=\small},
  usecase/.style={ellipse,draw=cLine,thick,fill=cBlue,align=center,font=\footnotesize,inner sep=2pt,minimum height=8mm},
  proc/.style={circle,draw=cLine,thick,fill=cBlue,align=center,font=\footnotesize,inner sep=2pt},
  arr/.style={->,thick,draw=cLine},
  darr/.style={->,thick,draw=cLine,dashed},
}

% ---- helper: include a figure file only if it exists (keeps the build alive) ---
\newcommand{\OptionalGraphic}[2][width=\linewidth]{%
  \IfFileExists{#2}{\includegraphics[#1]{#2}}{\fbox{\parbox{0.8\linewidth}{\centering\small Image file \texttt{#2} not found}}}}

% ---- signature block used on certificate and abstract pages --------------------
\newcommand{\SignBlock}[2]{\begin{tabular}[t]{@{}c@{}}\rule{4.2cm}{0.4pt}\\[-1pt]\textbf{#1}\\#2\end{tabular}}

% ---- cover page (the first argument is the report type line) -------------------
\newcommand{\CoverPage}[1]{%
  \thispagestyle{empty}
  \begin{center}
    {\bfseries #1\par}
    \vspace{10pt}
    {\bfseries ``\ProjectTitleCaps''\par}
    \vspace{14pt}
    SUBMITTED TO THE SAVITRIBAI PHULE PUNE UNIVERSITY, PUNE\\
    IN THE PARTIAL FULFILLMENT OF THE REQUIREMENTS\\
    FOR THE AWARD OF THE DEGREE\\[8pt]
    OF\\[8pt]
    {\bfseries BACHELOR OF ENGINEERING (COMPUTER ENGINEERING)}\par
    \vspace{22pt}
    {\bfseries SUBMITTED BY}\par
    \vspace{8pt}
    \begin{tabular}{@{}ll@{}}
      \StudentOne & PRN No.: \PRNOne\\
      \StudentTwo & PRN No.: \PRNTwo\\
      \StudentThree & PRN No.: \PRNThree
    \end{tabular}\par
    \vspace{16pt}
    \IfFileExists{bsiotr_logo.png}{\includegraphics[height=2.6cm]{bsiotr_logo.png}\par}{}
    \vspace{10pt}
    {\bfseries DEPARTMENT OF COMPUTER ENGINEERING}\par
    \vspace{10pt}
    {\bfseries JSPM'S, BHIVARABAI SAWANT INSTITUTE OF TECHNOLOGY \& RESEARCH}\par
    \InstituteAddress\par
    \vspace{6pt}
    {\bfseries SAVITRIBAI PHULE PUNE UNIVERSITY}\par
    \vspace{6pt}
    {\bfseries \AcademicYear}
  \end{center}
  \newpage}
.
% Only standard packages that are on Overleaf by default are used.
% ============================================================================
\usepackage[T1]{fontenc}
\usepackage[utf8]{inputenc}
\usepackage{mathptmx}                       % Times-style text and maths
\usepackage[a4paper,left=3cm,right=2.5cm,top=2.5cm,bottom=2.6cm]{geometry}
\usepackage{setspace}
\usepackage{graphicx}
\usepackage{xcolor}
\usepackage{booktabs,array,tabularx,multirow,longtable}
\usepackage{amsmath,amssymb}
\usepackage{enumitem}
\usepackage[font=small,labelfont=bf]{caption}
\usepackage{float}
\usepackage{fancyhdr}
\usepackage{titlesec}
\usepackage{url}
\usepackage[hidelinks]{hyperref}
\usepackage{tikz}
\usetikzlibrary{arrows.meta,positioning,shapes.geometric,shapes.misc,calc,fit,backgrounds,decorations.pathreplacing}
\usepackage{pgfplots}
\pgfplotsset{compat=1.16}
\usepackage{pgfgantt}

% ---- colours (also readable in black and white) ----------------------------
\definecolor{cBlue}{HTML}{DCE8F5}
\definecolor{cGreen}{HTML}{DDEEDD}
\definecolor{cOrange}{HTML}{FBE5CC}
\definecolor{cGrey}{HTML}{EEEEEE}
\definecolor{cLine}{HTML}{333333}

% ---- page style: footer as in the department's report format -----------------
\newcommand{\FooterText}{BSIOTR, Department of Computer Engineering, \AcademicYearShort}
\pagestyle{fancy}
\fancyhf{}
\renewcommand{\headrulewidth}{0pt}
\fancyfoot[C]{\small \FooterText\\[2pt]\thepage}
\fancypagestyle{plain}{\fancyhf{}\renewcommand{\headrulewidth}{0pt}\fancyfoot[C]{\small \FooterText\\[2pt]\thepage}}

% ---- headings ------------------------------------------------------------------
\makeatletter
\@ifundefined{chapter}{}{%
  \titleformat{\chapter}[display]{\centering\bfseries\Large}{CHAPTER \thechapter}{8pt}{\Large\MakeUppercase}
  \titlespacing*{\chapter}{0pt}{-10pt}{22pt}}
\makeatother
\titleformat{\section}{\bfseries\large}{\thesection}{0.6em}{}
\titleformat{\subsection}{\bfseries\normalsize}{\thesubsection}{0.6em}{}
\setlength{\parindent}{0pt}
\setlength{\parskip}{6pt}
\setlist{itemsep=2pt,topsep=3pt}
\renewcommand{\arraystretch}{1.18}
\setlength{\emergencystretch}{3em}

% ---- TikZ styles shared by all diagrams ----------------------------------------
\tikzset{
  >=Stealth,
  box/.style={rectangle,rounded corners=2pt,draw=cLine,thick,fill=cBlue,align=center,
              minimum height=9mm,inner sep=4pt,font=\small},
  boxg/.style={box,fill=cGreen},
  boxo/.style={box,fill=cOrange},
  boxw/.style={box,fill=white},
  designed/.style={box,fill=white,dashed},
  store/.style={rectangle,draw=cLine,thick,fill=cGrey,align=center,minimum height=8mm,inner sep=4pt,font=\small},
  actor/.style={rectangle,draw=cLine,thick,fill=cOrange,align=center,minimum height=8mm,inner sep=4pt,font=\small},
  usecase/.style={ellipse,draw=cLine,thick,fill=cBlue,align=center,font=\footnotesize,inner sep=2pt,minimum height=8mm},
  proc/.style={circle,draw=cLine,thick,fill=cBlue,align=center,font=\footnotesize,inner sep=2pt},
  arr/.style={->,thick,draw=cLine},
  darr/.style={->,thick,draw=cLine,dashed},
}

% ---- helper: include a figure file only if it exists (keeps the build alive) ---
\newcommand{\OptionalGraphic}[2][width=\linewidth]{%
  \IfFileExists{#2}{\includegraphics[#1]{#2}}{\fbox{\parbox{0.8\linewidth}{\centering\small Image file \texttt{#2} not found}}}}

% ---- signature block used on certificate and abstract pages --------------------
\newcommand{\SignBlock}[2]{\begin{tabular}[t]{@{}c@{}}\rule{4.2cm}{0.4pt}\\[-1pt]\textbf{#1}\\#2\end{tabular}}

% ---- cover page (the first argument is the report type line) -------------------
\newcommand{\CoverPage}[1]{%
  \thispagestyle{empty}
  \begin{center}
    {\bfseries #1\par}
    \vspace{10pt}
    {\bfseries ``\ProjectTitleCaps''\par}
    \vspace{14pt}
    SUBMITTED TO THE SAVITRIBAI PHULE PUNE UNIVERSITY, PUNE\\
    IN THE PARTIAL FULFILLMENT OF THE REQUIREMENTS\\
    FOR THE AWARD OF THE DEGREE\\[8pt]
    OF\\[8pt]
    {\bfseries BACHELOR OF ENGINEERING (COMPUTER ENGINEERING)}\par
    \vspace{22pt}
    {\bfseries SUBMITTED BY}\par
    \vspace{8pt}
    \begin{tabular}{@{}ll@{}}
      \StudentOne & PRN No.: \PRNOne\\
      \StudentTwo & PRN No.: \PRNTwo\\
      \StudentThree & PRN No.: \PRNThree
    \end{tabular}\par
    \vspace{16pt}
    \IfFileExists{bsiotr_logo.png}{\includegraphics[height=2.6cm]{bsiotr_logo.png}\par}{}
    \vspace{10pt}
    {\bfseries DEPARTMENT OF COMPUTER ENGINEERING}\par
    \vspace{10pt}
    {\bfseries JSPM'S, BHIVARABAI SAWANT INSTITUTE OF TECHNOLOGY \& RESEARCH}\par
    \InstituteAddress\par
    \vspace{6pt}
    {\bfseries SAVITRIBAI PHULE PUNE UNIVERSITY}\par
    \vspace{6pt}
    {\bfseries \AcademicYear}
  \end{center}
  \newpage}
.
% Only standard packages that are on Overleaf by default are used.
% ============================================================================
\usepackage[T1]{fontenc}
\usepackage[utf8]{inputenc}
\usepackage{mathptmx}                       % Times-style text and maths
\usepackage[a4paper,left=3cm,right=2.5cm,top=2.5cm,bottom=2.6cm]{geometry}
\usepackage{setspace}
\usepackage{graphicx}
\usepackage{xcolor}
\usepackage{booktabs,array,tabularx,multirow,longtable}
\usepackage{amsmath,amssymb}
\usepackage{enumitem}
\usepackage[font=small,labelfont=bf]{caption}
\usepackage{float}
\usepackage{fancyhdr}
\usepackage{titlesec}
\usepackage{url}
\usepackage[hidelinks]{hyperref}
\usepackage{tikz}
\usetikzlibrary{arrows.meta,positioning,shapes.geometric,shapes.misc,calc,fit,backgrounds,decorations.pathreplacing}
\usepackage{pgfplots}
\pgfplotsset{compat=1.16}
\usepackage{pgfgantt}

% ---- colours (also readable in black and white) ----------------------------
\definecolor{cBlue}{HTML}{DCE8F5}
\definecolor{cGreen}{HTML}{DDEEDD}
\definecolor{cOrange}{HTML}{FBE5CC}
\definecolor{cGrey}{HTML}{EEEEEE}
\definecolor{cLine}{HTML}{333333}

% ---- page style: footer as in the department's report format -----------------
\newcommand{\FooterText}{BSIOTR, Department of Computer Engineering, \AcademicYearShort}
\pagestyle{fancy}
\fancyhf{}
\renewcommand{\headrulewidth}{0pt}
\fancyfoot[C]{\small \FooterText\\[2pt]\thepage}
\fancypagestyle{plain}{\fancyhf{}\renewcommand{\headrulewidth}{0pt}\fancyfoot[C]{\small \FooterText\\[2pt]\thepage}}

% ---- headings ------------------------------------------------------------------
\makeatletter
\@ifundefined{chapter}{}{%
  \titleformat{\chapter}[display]{\centering\bfseries\Large}{CHAPTER \thechapter}{8pt}{\Large\MakeUppercase}
  \titlespacing*{\chapter}{0pt}{-10pt}{22pt}}
\makeatother
\titleformat{\section}{\bfseries\large}{\thesection}{0.6em}{}
\titleformat{\subsection}{\bfseries\normalsize}{\thesubsection}{0.6em}{}
\setlength{\parindent}{0pt}
\setlength{\parskip}{6pt}
\setlist{itemsep=2pt,topsep=3pt}
\renewcommand{\arraystretch}{1.18}
\setlength{\emergencystretch}{3em}

% ---- TikZ styles shared by all diagrams ----------------------------------------
\tikzset{
  >=Stealth,
  box/.style={rectangle,rounded corners=2pt,draw=cLine,thick,fill=cBlue,align=center,
              minimum height=9mm,inner sep=4pt,font=\small},
  boxg/.style={box,fill=cGreen},
  boxo/.style={box,fill=cOrange},
  boxw/.style={box,fill=white},
  designed/.style={box,fill=white,dashed},
  store/.style={rectangle,draw=cLine,thick,fill=cGrey,align=center,minimum height=8mm,inner sep=4pt,font=\small},
  actor/.style={rectangle,draw=cLine,thick,fill=cOrange,align=center,minimum height=8mm,inner sep=4pt,font=\small},
  usecase/.style={ellipse,draw=cLine,thick,fill=cBlue,align=center,font=\footnotesize,inner sep=2pt,minimum height=8mm},
  proc/.style={circle,draw=cLine,thick,fill=cBlue,align=center,font=\footnotesize,inner sep=2pt},
  arr/.style={->,thick,draw=cLine},
  darr/.style={->,thick,draw=cLine,dashed},
}

% ---- helper: include a figure file only if it exists (keeps the build alive) ---
\newcommand{\OptionalGraphic}[2][width=\linewidth]{%
  \IfFileExists{#2}{\includegraphics[#1]{#2}}{\fbox{\parbox{0.8\linewidth}{\centering\small Image file \texttt{#2} not found}}}}

% ---- signature block used on certificate and abstract pages --------------------
\newcommand{\SignBlock}[2]{\begin{tabular}[t]{@{}c@{}}\rule{4.2cm}{0.4pt}\\[-1pt]\textbf{#1}\\#2\end{tabular}}

% ---- cover page (the first argument is the report type line) -------------------
\newcommand{\CoverPage}[1]{%
  \thispagestyle{empty}
  \begin{center}
    {\bfseries #1\par}
    \vspace{10pt}
    {\bfseries ``\ProjectTitleCaps''\par}
    \vspace{14pt}
    SUBMITTED TO THE SAVITRIBAI PHULE PUNE UNIVERSITY, PUNE\\
    IN THE PARTIAL FULFILLMENT OF THE REQUIREMENTS\\
    FOR THE AWARD OF THE DEGREE\\[8pt]
    OF\\[8pt]
    {\bfseries BACHELOR OF ENGINEERING (COMPUTER ENGINEERING)}\par
    \vspace{22pt}
    {\bfseries SUBMITTED BY}\par
    \vspace{8pt}
    \begin{tabular}{@{}ll@{}}
      \StudentOne & PRN No.: \PRNOne\\
      \StudentTwo & PRN No.: \PRNTwo\\
      \StudentThree & PRN No.: \PRNThree
    \end{tabular}\par
    \vspace{16pt}
    \IfFileExists{bsiotr_logo.png}{\includegraphics[height=2.6cm]{bsiotr_logo.png}\par}{}
    \vspace{10pt}
    {\bfseries DEPARTMENT OF COMPUTER ENGINEERING}\par
    \vspace{10pt}
    {\bfseries JSPM'S, BHIVARABAI SAWANT INSTITUTE OF TECHNOLOGY \& RESEARCH}\par
    \InstituteAddress\par
    \vspace{6pt}
    {\bfseries SAVITRIBAI PHULE PUNE UNIVERSITY}\par
    \vspace{6pt}
    {\bfseries \AcademicYear}
  \end{center}
  \newpage}

\renewcommand{\thesubsection}{\thesection.\alph{subsection}}
\titleformat{\section}{\bfseries\large\MakeUppercase}{\thesection.}{0.6em}{}
\begin{document}

\CoverPage{A SYNOPSIS REPORT ON}

% ------------------------------------------------------------------- contents
\thispagestyle{empty}
\begin{center}{\bfseries\Large\underline{CONTENTS}}\end{center}
\vspace{6pt}
\renewcommand{\contentsname}{}
{\small\tableofcontents}
\newpage
\setcounter{page}{1}
\onehalfspacing

% ---------------------------------------------------------------- 1
\section{Introduction}
Images produced by AI tools such as Stable Diffusion, Midjourney, FLUX, the ChatGPT image model and the Gemini image model now look like photographs \cite{rombach}. Face-swapped video is common \cite{rossler}, and a figure on an invoice can be changed in a few minutes. Looking at such a file by eye is no longer a reliable way to judge it.

Research on detection is organised by medium: one line of work studies images, another face videos, another document forgery. The detectors are usually tuned to the generators they were trained on, so a detector trained on one generator can fall close to chance on the next \cite{ojha,wang}. Datasets can leak the answer through the file format \cite{grommelt}, and compression removes much of the evidence \cite{corvi}.

This project builds one web-based system that accepts an image, a video or a document, sends it to a specialised analyser and turns the evidence into a calibrated Trust Score from 0 to 100, a verdict and a plain-language explanation. The work is done in three phases (image, video, document). The image analyser is complete and evaluated; the video and document analysers are designed and will be built next.

% ---------------------------------------------------------------- 2
\section{Motivation / Need of the Project}
\begin{itemize}
  \item \textbf{Harm from fake content.} Synthetic images spread false news, fake documents are used in fraud, and fake pictures of real people are used to harass them.
  \item \textbf{No tool for ordinary users.} Benchmarks such as ForensicHub \cite{forensichub} bring many detectors together, but a benchmark is not a tool that a student, a journalist or a bank clerk can use.
  \item \textbf{New generators appear all the time.} A detector is only useful if it keeps working on tools it was not trained on, such as the latest ChatGPT and Gemini image models. This is the main research question of the project.
  \item \textbf{Uncertainty must be shown.} Calling a real photograph ``AI-generated'' can harm its author, so the system must be able to say ``a person should check this''.
  \item \textbf{Low cost.} The system must run on a CPU inside the memory limit of free hosting, so that anyone can use it without special hardware.
\end{itemize}

% ---------------------------------------------------------------- 3
\section{Problem Definition}
To design and develop a web-based, explainable system that checks whether an image, a video or a document is authentic, AI-generated or tampered with. The system routes each file to a specialised analyser and reports a calibrated Trust Score, a verdict with a human-review band and the evidence behind it. Detection must hold for generators that were not seen in training.

\textbf{Objectives}
\begin{itemize}
  \item To detect AI-generated images, including images from unseen generators, with an unseen-generator AUC of at least 0.80 and better than an EfficientNet baseline \cite{effnet}.
  \item To detect face-manipulated videos (FaceForensics++ and Celeb-DF) and to locate tampered regions in documents.
  \item To convert every result into one Trust Score with a calibration error of at most 0.05 and a ``Needs Human Review'' band.
  \item To explain each verdict with a heatmap and sentences that are built only from measured evidence.
  \item To keep the system on a CPU with under 2.7 GB of memory, with no uploads stored.
\end{itemize}

% ---------------------------------------------------------------- 4
\section{Existing System}
\subsection{Methodology (Architecture)}
A typical existing detector is built in three steps: collect real images and fake images from one or a few generators, train a classifier on them, and test it on images made by the same generators. Wang et al. trained a CNN with blur and JPEG augmentation \cite{wang}. Ojha et al. trained a linear classifier on frozen CLIP features and showed that it transfers better to unseen generators \cite{ojha}. Public benchmarks such as GenImage \cite{genimage} and OpenFake \cite{openfake} provide images from many generators. For video, FaceForensics++ \cite{rossler} and Celeb-DF \cite{celebdf} are the standard sets, and for documents DocTamper \cite{doctamper} is used.

\begin{figure}[htbp]
  \centering
  % Existing approach: a detector trained on one source and the problems that follow
\resizebox{\linewidth}{!}{%
\begin{tikzpicture}
\node[box,text width=3.0cm] (d) at (0,0) {Real photographs and images from \emph{one} generator};
\node[box,text width=3.0cm] (m) at (4.4,0) {CNN or frequency-based classifier trained end to end};
\node[box,text width=3.0cm] (t1) at (8.8,1.5) {Test on the same generator};
\node[box,text width=3.0cm] (t2) at (8.8,-1.5) {Test on a new generator or a compressed file};
\node[boxg,text width=3.0cm] (r1) at (13.2,1.5) {High accuracy};
\node[boxo,text width=3.0cm] (r2) at (13.2,-1.5) {Accuracy can fall close to chance};
\draw[arr] (d) -- (m);
\draw[arr] (m.east) -- ++(0.7,0) |- (t1.west);
\draw[arr] (m.east) -- ++(0.7,0) |- (t2.west);
\draw[arr] (t1) -- (r1);
\draw[arr] (t2) -- (r2);
\node[font=\footnotesize,text width=5.4cm,align=center] at (4.4,-2.6) {Further weaknesses: dataset shortcuts such as file format, no calibrated confidence, no explanation};
\end{tikzpicture}}

  \caption{Typical existing approach: training on one source and the problems that follow.}
\end{figure}

\subsection{Advantages and Disadvantages}
\textbf{Advantages}
\begin{itemize}
  \item The methods are well studied and public benchmarks exist for images, face video and documents.
  \item Detectors built on CLIP features generalise better than detectors trained from scratch \cite{ojha}.
\end{itemize}
\textbf{Disadvantages}
\begin{itemize}
  \item Many detectors are tied to the generators they saw in training and can fall close to chance on new ones.
  \item Some datasets leak the label through file format \cite{grommelt}, and compression or resizing weakens most detectors \cite{corvi}.
  \item Each system handles one medium only, and most give a bare label without a calibrated confidence or an explanation.
\end{itemize}

% ---------------------------------------------------------------- 5
\section{Proposed System}
\subsection{Methodology (Architecture)}
The proposed system has a layered design (Fig.~\ref{fig:arch}). A Streamlit web interface receives the upload; a validation layer reads the true file type from the bytes and decodes the file; a router calls the image, video or document analyser. All analysers return the same kind of evidence record (a score, a visual, numeric cues and flags). A scoring layer calibrates the record and an explanation layer writes the sentences. Because the analysers share only this interface, one of them can be replaced without touching the rest of the application.

\begin{figure}[htbp]
  \centering
  % Layered architecture: solid = implemented, dashed = designed
\resizebox{\linewidth}{!}{%
\begin{tikzpicture}
\node[box,text width=3.8cm] (up) at (1.9,0) {User upload\\image, video or document};
\node[box,text width=5.4cm] (val) at (7.6,0) {Validation and routing\\file type read from the bytes,\\size limit, safe decoding};
\draw[arr] (up) -- (val);
% bus to the three analysers
\draw[thick,cLine] (val.south) -- (7.6,-1.25);
\draw[thick,cLine] (1.9,-1.25) -- (13.3,-1.25);
\node[boxg,text width=4.1cm] (img) at (1.9,-3.0) {\textbf{Image analyser}\\(implemented)\\native-resolution patches,\\frozen CLIP + trained head,\\patch heatmap};
\node[designed,text width=4.1cm] (vid) at (7.6,-3.0) {\textbf{Video analyser}\\(designed)\\sampled frames, face crops,\\frame scores, timeline};
\node[designed,text width=4.1cm] (doc) at (13.3,-3.0) {\textbf{Document analyser}\\(designed)\\OCR, ELA / SRM / DCT views,\\dual-stream U-Net mask,\\rule checks};
\draw[arr] (1.9,-1.25) -- (img.north);
\draw[arr] (7.6,-1.25) -- (vid.north);
\draw[arr] (13.3,-1.25) -- (doc.north);
% evidence bus
\draw[thick,cLine] (img.south) -- (1.9,-5.0);
\draw[thick,cLine] (vid.south) -- (7.6,-5.0);
\draw[thick,cLine] (doc.south) -- (13.3,-5.0);
\draw[thick,cLine] (1.9,-5.0) -- (13.3,-5.0);
\node[box,text width=6.4cm] (ev) at (7.6,-6.1) {Evidence records\\score, map or timeline, numeric cues, flags};
\draw[arr] (7.6,-5.0) -- (ev.north);
\node[box,text width=6.4cm] (cal) at (7.6,-7.7) {Calibration and fusion\\$p=\sigma(a\,s+c)$ (Platt scaling)};
\draw[arr] (ev) -- (cal);
\node[boxo,text width=8.0cm] (out) at (7.6,-9.3) {Trust Score 0--100, verdict band,\\explanation sentences, downloadable report};
\draw[arr] (cal) -- (out);
\end{tikzpicture}}

  \caption{Layered architecture. Solid boxes are implemented; dashed boxes are designed.}\label{fig:arch}
\end{figure}

Training and use are kept apart (Fig.~\ref{fig:train}). The heavy work, such as turning hundreds of thousands of patches into features, runs once on a free cloud GPU. The application loads the frozen encoder and a trained head of a few kilobytes and runs on a CPU.

\begin{figure}[htbp]
  \centering
  % Training (offline) and inference (online) workflow
\resizebox{\linewidth}{!}{%
\begin{tikzpicture}
\node[font=\small\bfseries] at (7.2,1.0) {Offline: free cloud GPU and laptop};
\node[box,text width=2.9cm] (a1) at (0,0) {Public datasets\\(OpenFake, DiffSeg30k)};
\node[box,text width=2.9cm] (a2) at (3.7,0) {Patches at native size + degraded copies};
\node[box,text width=2.9cm] (a3) at (7.4,0) {Frozen CLIP features, stored once};
\node[box,text width=2.9cm] (a4) at (11.1,0) {Fit small head, choose $C$ on validation};
\node[box,text width=2.9cm] (a5) at (14.8,0) {Platt scaling and verdict limits on validation};
\draw[arr] (a1) -- (a2);
\draw[arr] (a2) -- (a3);
\draw[arr] (a3) -- (a4);
\draw[arr] (a4) -- (a5);
\node[boxw,text width=3.6cm] (t) at (14.8,-2.2) {Held-out test split scored once};
\draw[arr] (a5) -- (t);
\node[store,text width=3.6cm] (w) at (9.2,-2.2) {Weight files\\(a few kilobytes)};
\draw[arr] (t) -- (w);
\node[font=\small\bfseries] at (7.2,-3.6) {Online: CPU only};
\node[boxg,text width=2.9cm] (b1) at (0,-5.0) {Upload and validate};
\node[boxg,text width=2.9cm] (b2) at (3.7,-5.0) {Ten patches, JPEG q95};
\node[boxg,text width=2.9cm] (b3) at (7.4,-5.0) {CLIP encoder + loaded head};
\node[boxg,text width=2.9cm] (b4) at (11.1,-5.0) {Calibrated Trust Score};
\node[boxo,text width=2.9cm] (b5) at (14.8,-5.0) {Verdict, heatmap, explanation};
\draw[arr] (b1) -- (b2);
\draw[arr] (b2) -- (b3);
\draw[arr] (b3) -- (b4);
\draw[arr] (b4) -- (b5);
\draw[darr] (w.south) -- ++(0,-1.3) -| (b3.north);
\end{tikzpicture}}

  \caption{Offline training workflow and online use of the trained model.}\label{fig:train}
\end{figure}

\subsection{Project Modules}
The basic modules of the project are:
\begin{enumerate}[label=\arabic*)]
  \item Input validation and routing
  \item Image analyser (implemented)
  \item Video analyser (designed)
  \item Document analyser (designed)
  \item Scoring, calibration and explanation
  \item Web interface and report
\end{enumerate}

\textbf{1) Input validation and routing.} The file type is read from the bytes and not from the file name; size limits are applied; damaged files are refused with a clear message. Uploads are processed in memory and are not stored.

\textbf{2) Image analyser.} The image is not resized, because resizing removes fine traces. Ten $224\times224$ patches (the centre and a $3\times3$ grid) are cut at native resolution and re-saved as JPEG with quality 95, so that real and generated inputs share one compression history. Each patch is encoded by a frozen CLIP ViT-L/14 model and scored by a linear head; the image score is the mean of the patch scores, and the patch scores form the heatmap (Fig.~\ref{fig:pipe}).

\begin{figure}[htbp]
  \centering
  % Image analyser pipeline (two rows)
\resizebox{\linewidth}{!}{%
\begin{tikzpicture}
\node[box,text width=3.0cm] (b1) at (0,0) {\textbf{1.} Input image\\(not resized)};
\node[box,text width=3.0cm] (b2) at (4.0,0) {\textbf{2.} Ten patches\\224$\times$224: centre\\+ 3$\times$3 grid};
\node[box,text width=3.0cm] (b3) at (8.0,0) {\textbf{3.} Re-save each patch\\as JPEG, quality 95};
\node[box,text width=3.0cm] (b4) at (12.0,0) {\textbf{4.} Frozen CLIP\\ViT-L/14 encoder,\\$z=f(x)/\lVert f(x)\rVert$};
\draw[arr] (b1) -- (b2);
\draw[arr] (b2) -- (b3);
\draw[arr] (b3) -- (b4);
\node[boxg,text width=3.0cm] (b5) at (0,-3.0) {\textbf{5.} Linear head\\$s=\mathbf{w}^{\mathsf T}z+b$\\per patch};
\node[boxg,text width=3.0cm] (b6) at (4.0,-3.0) {\textbf{6.} Image score\\= mean of the ten\\patch scores};
\node[boxg,text width=3.0cm] (b7) at (8.0,-3.0) {\textbf{7.} Platt scaling\\$p=\sigma(a\,s+c)$};
\node[boxo,text width=3.0cm] (b8) at (12.0,-3.0) {\textbf{8.} Trust Score\\$T=\mathrm{round}(100(1-p))$\\and verdict band};
\draw[arr] (b4.south) -- ++(0,-0.75) -| (b5.north);
\draw[arr] (b5) -- (b6);
\draw[arr] (b6) -- (b7);
\draw[arr] (b7) -- (b8);
\node[font=\footnotesize,text width=3.0cm,align=center] at (4.0,-4.6) {patch scores also form the $3\times3$ heatmap};
\draw[darr] (4.0,-4.25) -- (4.0,-3.55);
\end{tikzpicture}}

  \caption{Processing steps of the image analyser.}\label{fig:pipe}
\end{figure}

\textbf{3) Video analyser (designed).} About one frame per second is sampled (at most 32). MTCNN \cite{mtcnn} finds and aligns the faces; every face crop is scored, first by the image analyser unchanged and then by a head trained on FaceForensics++. The frame scores are averaged over the top-scoring share of frames, so a short manipulated segment is not hidden by clean frames, and they also give a timeline.

\textbf{4) Document analyser (designed).} Pages are rendered and read by OCR \cite{tesseract}. Error-level analysis \cite{krawetz}, noise residuals from the spatial rich model \cite{fridrich} and DCT features are stacked as a second input stream next to the RGB page of a dual-stream U-Net \cite{unet,zhou}. Rule checks (line items against the total, dates, font heights, PDF history) add flags that are auxiliary evidence, never proof.

\textbf{5) Scoring, calibration and explanation.} A raw score $s$ becomes a probability $p=\sigma(a\,s+c)$ by Platt scaling \cite{platt}, fitted on validation data. The Trust Score is $T=\mathrm{round}\bigl(100(1-p)\bigr)$. Two limits chosen on validation data divide $T$ into \emph{Likely Authentic} ($T\ge 54$), \emph{Needs Human Review} and \emph{Likely AI-generated / Manipulated} ($T\le 17$). The sentences of the explanation come only from measured evidence.

\textbf{6) Web interface and report.} The user uploads a file and sees the Trust Score, the verdict, the heatmap, the explanation, the file metadata (shown apart from the score) and a downloadable report.

\subsection{Algorithms Used}
\textbf{1) CLIP image encoder (frozen).} The ViT-L/14 image encoder of CLIP \cite{radford} turns a patch into a 768-number feature vector. It is not trained; only the head on top of it is.

\textbf{2) Logistic regression head.} A linear head $s=\mathbf{w}^{\mathsf T}z+b$ on the normalised feature $z=f(x)/\lVert f(x)\rVert$, trained with L2 regularisation on features of real photographs and images from about sixty generators, including edited-image patches \cite{diffseg}. Every training image also gives a randomly resized, blurred and JPEG-compressed copy, an idea taken from blur-and-JPEG augmentation \cite{wang}.

\textbf{3) Patch averaging.} The image score is the mean of the ten patch scores, so that a few unusual patches do not decide the result.

\textbf{4) Platt scaling and verdict limits.} Calibration \cite{guo,platt} makes the probability behave like a real probability; the verdict limits are chosen on validation data so that at most 3\% of real images are called AI-generated and at most 10\% of AI images are called authentic \cite{geifman}.

\textbf{5) Video and document methods (planned).} MTCNN \cite{mtcnn}, top-fraction aggregation, ELA \cite{krawetz}, SRM \cite{fridrich}, dual-stream U-Net \cite{unet} with binary cross-entropy plus a Dice term \cite{vnet}.

\textbf{6) Evaluation.} ROC-AUC \cite{fawcett}, the share of AI images caught, the share of real images accepted, 95\% bootstrap intervals \cite{efron} and the expected calibration error \cite{naeini}.

\subsection{Advantages and Disadvantages}
\textbf{Advantages}\\
The image analyser was tested on 20 generators that were not in the training data (1,200 real and 2,192 AI images).

\begin{table}[htbp]
\centering\small
\caption{Image analyser on the held-out test split (clean files).}
\begin{tabular}{lcc}
\toprule
Test images & ROC-AUC & AI images caught\\
\midrule
All 20 unseen generators & 0.982 & 94\%\\
ChatGPT (GPT-Image-1.5) & 0.972 & 89\%\\
Gemini (Nano Banana Pro) & 0.985 & 94\%\\
Proof model, ChatGPT & 0.957 & 84\%\\
Proof model, Gemini & 0.984 & 98\%\\
\bottomrule
\end{tabular}
\end{table}

The ``proof model'' was trained without any OpenAI or Google image, so its result is a true unseen-source test. An EfficientNet-B0 baseline trained on the same images reaches an AUC of 0.831 on this test, against 0.982 for the proposed analyser. The model runs on a CPU, gives a calibrated score and an explanation, and keeps uploads private.

\textbf{Disadvantages}
\begin{itemize}
  \item The head is linear and the heatmap is coarse ($3\times3$ patches).
  \item Heavy compression and resizing weaken detection, and a small edited region can be diluted when patch scores are averaged.
  \item Only the image analyser is built and evaluated so far; the video and document analysers are designed.
  \item The Trust Score is a screening aid and not forensic proof.
\end{itemize}

% ---------------------------------------------------------------- 6
\section{Applications}
\begin{itemize}
  \item Fact-checking and newsrooms: screening pictures that arrive from the public.
  \item Social media and content moderation: flagging likely AI images for a human reviewer.
  \item Banks, insurers and KYC teams: first-level screening of submitted documents (planned for the document analyser).
  \item Education and awareness: showing students why a picture looks synthetic, through the heatmap and explanation.
  \item Personal use: checking a picture received in a messaging group.
\end{itemize}

% ---------------------------------------------------------------- 7
\section{Conclusion and Future Scope}
\textbf{Conclusion.} A frozen CLIP encoder with a small trained head, a degraded-copy training set and a calibrated, explainable scoring layer detects most AI images from unseen generators, and a layered design lets the same Trust Score serve images, videos and documents.

\textbf{Future scope.}
\begin{itemize}
  \item Build and evaluate the video analyser (Phase 2) on FaceForensics++ and Celeb-DF.
  \item Build the document analyser (Phase 3) with a dual-stream U-Net and rule checks, and test it on human-made forgeries and DocTamper.
  \item Improve robustness to compression and resizing, and finer, pixel-level maps for edited regions.
  \item Add more recent generators to the unseen-source tests as they appear.
  \item Offer a separate front end that calls the analysers through an API.
\end{itemize}

% ---------------------------------------------------------------- 8
\section{References}
\begingroup
\small
\begin{thebibliography}{99}
% Reference list: the same 32 references, in the same order, as the research paper.
\bibitem{rombach} R. Rombach \emph{et al.}, ``High-resolution image synthesis with latent diffusion models,'' in \emph{Proc. CVPR}, 2022, pp. 10684--10695.
\bibitem{rossler} A. R\"ossler \emph{et al.}, ``FaceForensics++: Learning to detect manipulated facial images,'' in \emph{Proc. ICCV}, 2019, pp. 1--11.
\bibitem{ojha} U. Ojha, Y. Li, and Y. J. Lee, ``Towards universal fake image detectors that generalize across generative models,'' in \emph{Proc. CVPR}, 2023, pp. 24480--24489.
\bibitem{wang} S.-Y. Wang \emph{et al.}, ``CNN-generated images are surprisingly easy to spot\ldots{} for now,'' in \emph{Proc. CVPR}, 2020, pp. 8695--8704.
\bibitem{grommelt} P. Grommelt, L. Weiss, F.-J. Pfreundt, and J. Keuper, ``Fake or JPEG? Revealing common biases in generated image detection datasets,'' arXiv:2403.17608, 2024.
\bibitem{corvi} R. Corvi \emph{et al.}, ``On the detection of synthetic images generated by diffusion models,'' in \emph{Proc. ICASSP}, 2023, arXiv:2211.00680.
\bibitem{forensichub} B. Du \emph{et al.}, ``ForensicHub: A unified benchmark \& codebase for all-domain fake image detection and localization,'' in \emph{NeurIPS Datasets and Benchmarks}, 2025, arXiv:2505.11003.
\bibitem{radford} A. Radford \emph{et al.}, ``Learning transferable visual models from natural language supervision,'' in \emph{Proc. ICML}, 2021, pp. 8748--8763.
\bibitem{genimage} M. Zhu \emph{et al.}, ``GenImage: A million-scale benchmark for detecting AI-generated image,'' in \emph{NeurIPS Datasets and Benchmarks}, 2023, arXiv:2306.08571.
\bibitem{openfake} V. Livernoche \emph{et al.}, ``OpenFake: An open dataset and platform toward real-world deepfake detection,'' arXiv:2509.09495, 2025.
\bibitem{celebdf} Y. Li, X. Yang, P. Sun, H. Qi, and S. Lyu, ``Celeb-DF: A large-scale challenging dataset for DeepFake forensics,'' in \emph{Proc. CVPR}, 2020, pp. 3207--3216.
\bibitem{mtcnn} K. Zhang, Z. Zhang, Z. Li, and Y. Qiao, ``Joint face detection and alignment using multitask cascaded convolutional networks,'' \emph{IEEE Signal Process. Lett.}, vol. 23, no. 10, pp. 1499--1503, 2016.
\bibitem{doctamper} C. Qu \emph{et al.}, ``Towards robust tampered text detection in document image: New dataset and new solution,'' in \emph{Proc. CVPR}, 2023, pp. 5937--5946.
\bibitem{zhou} P. Zhou, X. Han, V. I. Morariu, and L. S. Davis, ``Learning rich features for image manipulation detection,'' in \emph{Proc. CVPR}, 2018, pp. 1053--1061.
\bibitem{fridrich} J. Fridrich and J. Kodovsky, ``Rich models for steganalysis of digital images,'' \emph{IEEE Trans. Inf. Forensics Security}, vol. 7, no. 3, pp. 868--882, 2012.
\bibitem{krawetz} N. Krawetz, ``A picture's worth: Digital image analysis and forensics,'' Black Hat Briefings, 2007.
\bibitem{unet} O. Ronneberger, P. Fischer, and T. Brox, ``U-Net: Convolutional networks for biomedical image segmentation,'' in \emph{Proc. MICCAI}, 2015, pp. 234--241.
\bibitem{vnet} F. Milletari, N. Navab, and S.-A. Ahmadi, ``V-Net: Fully convolutional neural networks for volumetric medical image segmentation,'' in \emph{Proc. 3DV}, 2016, pp. 565--571.
\bibitem{guo} C. Guo, G. Pleiss, Y. Sun, and K. Q. Weinberger, ``On calibration of modern neural networks,'' in \emph{Proc. ICML}, 2017, pp. 1321--1330.
\bibitem{platt} J. C. Platt, ``Probabilistic outputs for support vector machines and comparisons to regularized likelihood methods,'' in \emph{Advances in Large Margin Classifiers}, MIT Press, 1999, pp. 61--74.
\bibitem{geifman} Y. Geifman and R. El-Yaniv, ``Selective classification for deep neural networks,'' in \emph{Proc. NeurIPS}, 2017.
\bibitem{effnet} M. Tan and Q. V. Le, ``EfficientNet: Rethinking model scaling for convolutional neural networks,'' in \emph{Proc. ICML}, 2019, pp. 6105--6114.
\bibitem{c2pa} Coalition for Content Provenance and Authenticity, ``C2PA technical specification,'' 2024. [Online]. Available: c2pa.org/specifications
\bibitem{pytorch} A. Paszke \emph{et al.}, ``PyTorch: An imperative style, high-performance deep learning library,'' in \emph{Proc. NeurIPS}, 2019, pp. 8024--8035.
\bibitem{sklearn} F. Pedregosa \emph{et al.}, ``Scikit-learn: Machine learning in Python,'' \emph{J. Mach. Learn. Res.}, vol. 12, pp. 2825--2830, 2011.
\bibitem{diffseg} H. Ci, Z. Peng, P. Yang, Y. Xuan, and M. Z. Shou, ``DiffSeg30k: A multi-turn diffusion editing benchmark for localized AIGC detection,'' arXiv:2511.19111, 2025.
\bibitem{li} L. Li and S. Lyu, ``Exposing DeepFake videos by detecting face warping artifacts,'' arXiv:1811.00656, 2018.
\bibitem{tesseract} R. Smith, ``An overview of the Tesseract OCR engine,'' in \emph{Proc. ICDAR}, 2007, pp. 629--633.
\bibitem{naeini} M. P. Naeini, G. F. Cooper, and M. Hauskrecht, ``Obtaining well calibrated probabilities using Bayesian binning into quantiles,'' in \emph{Proc. AAAI}, 2015, pp. 2901--2907.
\bibitem{gptedit} Y. Wang, S. Yang, B. Zhao, L. Zhang, Q. Liu, Y. Zhou, and C. Xie, ``GPT-Image-Edit-1.5M: A million-scale, GPT-generated image dataset,'' arXiv:2507.21033, 2025.
\bibitem{fawcett} T. Fawcett, ``An introduction to ROC analysis,'' \emph{Pattern Recognit. Lett.}, vol. 27, no. 8, pp. 861--874, 2006.
\bibitem{efron} B. Efron, ``Bootstrap methods: Another look at the jackknife,'' \emph{Ann. Statist.}, vol. 7, no. 1, pp. 1--26, 1979.

\end{thebibliography}
\endgroup

\end{document}
